\documentclass[10pt,letterpaper]{article}



\usepackage{fontspec}

\setmainfont{Latin Modern Roman}



\usepackage{amsmath,amssymb}

\usepackage[margin=0.8in]{geometry}

\usepackage{booktabs,longtable,array}

\usepackage{enumitem}

\usepackage[hidelinks]{hyperref}



\setlength{\parindent}{0pt}

\setlength{\parskip}{5pt}

\setlength{\emergencystretch}{3em}

\setlist[itemize]{leftmargin=1.5em,itemsep=2pt,topsep=4pt}

\setlist[enumerate]{leftmargin=1.7em,itemsep=5pt,topsep=5pt}

\renewcommand{\arraystretch}{1.15}



\title{STAT 1210: Assignment 3\\[4pt]\large Part III: Foundations of Inference}
\author{Fall 2026 Semester}
\date{}
\begin{document}
\maketitle
\begin{center}
\begin{tabular}{lr}
\toprule
\textbf{Question} & \textbf{Suggested marks} \\
\midrule
Question 1 & 20 \\
Question 2 & 20 \\
Question 3 & 20 \\
Question 4 & 20 \\
Question 5 & 20 \\
\midrule
\textbf{Total} & \textbf{100} \\
\bottomrule
\end{tabular}
\end{center}
\section*{Instructions}
\begin{itemize}
 \item Your work should show all steps clearly for each question.
 \item Your solutions must be handwritten on lined or letter-sized paper. Your name and student number must appear on the first page.
 \item Your name should also appear on every additional page, and all pages should be stapled together.
 \item R may be used to check numerical results, but the requested calculations, standardizations and statistical reasoning must be shown.
 \item Report numerical answers to at least three decimal places when appropriate.
 \item Programming resources, textbooks and software manuals may be consulted, but all submitted work must be your own. Solutions generated by generative AI tools are not permitted.
\end{itemize}
\newpage
% =======================================
% Question 1
% =======================================
\subsection*{Question 1}
A veterinary epidemiology team wants to study whether regular physical activity is associated with obesity among adult dogs in Prince Edward Island. Five veterinary clinics volunteered to participate. During a two-week period, owners arriving for appointments were invited to complete a questionnaire about their dog's average daily activity. A veterinarian then classified each participating dog as obese or not obese.
The results for the $600$ participating dogs are shown below.
\begin{center}
\begin{tabular}{lrrr}
\toprule
\textbf{Average daily activity} & \textbf{Obese} & \textbf{Not obese} & \textbf{Total} \\\\
\midrule
At least $60$ minutes & $18$ & $132$ & $150$ \\\\
Less than $60$ minutes & $126$ & $324$ & $450$ \\\\
\midrule
\textbf{Total} & $144$ & $456$ & $600$ \\\\
\bottomrule
\end{tabular}
\end{center}
\begin{enumerate}
 \item Identify the individuals, the intended population and the sample in this study. Identify the two main variables and classify each as categorical or quantitative.
 \item Calculate the proportion of obese dogs separately for the two activity groups. Describe the association between daily activity and obesity in the sample.
 \item Is this an observational study or a randomized experiment? Was random sampling used? Was random assignment used? Explain each answer.
 \item Identify two plausible confounding variables that could help explain the observed association. For each variable, explain how it could be related to both a dog's activity level and obesity status.
 \item Identify two ways in which the recruitment process could produce selection bias, undercoverage, voluntary-response bias or nonresponse bias. Then explain why the study cannot establish causation or confidently generalize its results to all adult dogs in Prince Edward Island.
\end{enumerate}
\newpage
% =======================================
% Question 2
% =======================================
\subsection*{Question 2}
A fisheries agency estimates that the probability that a tagged juvenile Atlantic salmon survives a particular downstream migration route is $p=0.42$. Suppose a random sample of $n=400$ tagged salmon is followed under similar environmental conditions. Let $\hat{p}$ be the proportion of salmon in the sample that survive the route.
Assume that the survival outcomes can be treated as independent and that the population is sufficiently large relative to the sample.
\begin{enumerate}
 \item Identify the population parameter and the sample statistic. Explain why the survival outcome for one salmon is categorical but $\hat{p}$ is quantitative.
 \item Calculate the mean and standard deviation of the sampling distribution of $\hat{p}$ using
 \[
 \mu_{\hat{p}}=p
 \qquad\text{and}\qquad
 \sigma_{\hat{p}}=\sqrt{\frac{p(1-p)}{n}}
 \]



 Interpret both quantities in context.



 \item Verify the expected-count conditions required for an approximately Normal sampling distribution of $\hat{p}$.



 \item Use the Normal approximation to calculate



 \[
 P(\hat{p}\leq0.38)
 \]



 Show the standardization calculation before using the Standard Normal table.



 \item Use the Normal approximation to calculate



 \[
 P(0.40\leq\hat{p}\leq0.46)
 \]



 Show both standardized boundary values.



 After calculating this probability, determine the sample size required to reduce the standard deviation of $\hat{p}$ to one-half of its value for $n=400$. Explain using the square-root relationship.

\end{enumerate}



\newpage



% =======================================

% Question 3

% =======================================



\subsection*{Question 3}



A veterinary hospital network wants to estimate the population mean recovery time following a routine anaesthetic procedure for adult dogs. The network performed this procedure on more than $1000$ adult dogs during the study period. A simple random sample of $64$ dogs has a sample mean recovery time of $47.8$ minutes.



Based on long-term monitoring of the same procedure, the population standard deviation is treated as known at $\sigma=12$ minutes.



\begin{enumerate}

 \item Define the population parameter $\mu$ in the context of this study. What statistic is being used as its point estimate?



 \item Assess the random-sampling, independence, $10\%$ and Normal-sampling-distribution conditions for using a one-sample $z$ confidence interval. State any assumption that cannot be checked directly from the information provided.



 \item Construct a $95\%$ confidence interval for the population mean recovery time. Show the standard error, critical value, margin of error and interval endpoints.



 \item Interpret the $95\%$ confidence interval in the context of the hospital network. Your interpretation should refer to the population mean rather than to individual recovery times.



 \item Construct a $99\%$ confidence interval using $z^*=2.576$. Compare its width with the width of the $95\%$ interval and explain why the widths differ.




\end{enumerate}



\newpage



% =======================================

% Question 4

% =======================================



\subsection*{Question 4}



Under an established sedation protocol, the population mean resting respiratory rate of adult harbour seals during recovery is reported to be $18$ breaths per minute. A marine veterinary team introduces a revised protocol. Because respiratory rates that are either too low or too high may be clinically important, the team wants to determine whether the population mean under the revised protocol differs from $18$ breaths per minute.



One respiratory-rate measurement is obtained from each of $25$ randomly selected adult harbour seals treated under the revised protocol.



\begin{center}

\begin{tabular}{rrrrr}

\toprule

$18$ & $22$ & $19$ & $23$ & $17$ \\\\

$21$ & $20$ & $23$ & $16$ & $23$ \\\\

$19$ & $18$ & $22$ & $21$ & $20$ \\\\

$24$ & $15$ & $23$ & $19$ & $20$ \\\\

$23$ & $18$ & $21$ & $17$ & $21$ \\\\

\bottomrule

\end{tabular}

\end{center}



Long-term monitoring indicates that respiratory rates under comparable conditions are approximately Normal with population standard deviation $\sigma=4$ breaths per minute. Use a significance level of $\alpha=0.05$.



\begin{enumerate}

 \item Define the population parameter $\mu$. State the null and alternative hypotheses, and explain why the alternative hypothesis is two-sided.



 \item Calculate the sample mean, the standard error and the one-sample $z$ test statistic. Interpret the sign and magnitude of the test statistic in context.



 \item Calculate the two-sided P-value using a Standard Normal table. Show how the probability in one tail is converted into a two-sided probability.



 Compare the P-value with $\alpha=0.05$, state the decision about $H_0$ and write a conclusion in context.



 \item Describe a Type I error and a Type II error in the context of this study. Which error might be more serious from an animal-safety perspective? Explain your reasoning.



 \item Construct a $95\%$ confidence interval for the population mean respiratory rate. Determine whether $18$ lies inside the interval and explain how this agrees with the test decision. Does statistical significance alone establish clinical importance? Explain briefly.

\end{enumerate}



% =======================================
% Question 5
% =======================================

\newpage

\subsection*{Question 5}

A veterinary diagnostic laboratory is evaluating a portable blood-glucose meter for use with canine blood samples. For each sample, the laboratory calculates

\[
\text{difference}
=
\text{portable-meter reading}
-
\text{reference-laboratory reading}
\]

A positive difference means that the portable meter tends to read higher than the reference method. A simple random sample of $400$ stored canine blood samples gives a mean difference of $\bar{x}=0.30$ mg/dL. Based on previous validation work, the population standard deviation of the differences is treated as known at $\sigma=2.4$ mg/dL.

The laboratory tests

\[
H_0:\mu=0
\qquad\text{against}\qquad
H_a:\mu\ne0
\]

at the significance level $\alpha=0.05$. For operational purposes, the laboratory considers a population mean bias with magnitude less than $1.0$ mg/dL to be too small to be practically important.

\begin{enumerate}

\item Define the population parameter $\mu$ in context. Calculate the standard error, the one-sample $z$ statistic and the two-sided P-value. State the decision about $H_0$ at $\alpha=0.05$.

\item Construct a $95\%$ confidence interval for $\mu$. Explain how the interval leads to the same decision as the two-sided test in Part 1.

\item Distinguish between statistical significance and practical importance for this study. Use the confidence interval and the laboratory's $1.0$ mg/dL threshold in your explanation. Identify one possible source of systematic error or bias that is not included in the interval's margin of error.

\item Construct a $99\%$ confidence interval using $z^*=2.576$. Compare it with the $95\%$ interval and explain why the corresponding test decision at $\alpha=0.01$ differs from the decision at $\alpha=0.05$.

\item For a future study, suppose the laboratory is concerned only with positive bias and plans to test $H_0:\mu=0$ against $H_a:\mu>0$. The laboratory wants $90\%$ power to detect a true mean bias of $\delta=0.6$ mg/dL using $\alpha=0.05$ and $\sigma=2.4$ mg/dL. Calculate the required sample size using

\[
n
=
\left(
\frac{(z_{1-\alpha}+z_{\text{power}})\sigma}{\delta}
\right)^2
\]

with $z_{1-\alpha}=1.645$ and $z_{\text{power}}=1.282$. Explain why the result must be rounded up.

\end{enumerate}



\end{document}

